\documentclass[a4paper]{article}
\input{../../../../../newpreamble.tex}
% \usepackage[nohead]{geometry}
% \title{Math 210A: Exam 2}
% \author{Lance Remigio}
\begin{document}
% \maketitle
\lhead{Math 210A: Midterm 1}
\chead{Lance Remigio}
\rhead{\thepage}
\begin{problem}
    Let \( \varphi: {G}_{1} \to {G}_{2} \) be a surjective group homomorphism. 
    \begin{enumerate}
        \item[(a)] Suppose that \( {G}_{1} = \Z / 39 \Z   \) and \( | {G}_{2} |  = 13 \). Determine with proof the kernel of \( \varphi \). 
            \begin{proof}
            Since \( \varphi  \) is a group homomorphism, we know by the first isomorphism theorem that 
            \begin{align*}
                {G}_{1} / \text{ker}(\varphi) \cong \varphi({G}_{1}) = {G}_{2} &\implies | {G}_{1} / \text{ker}(\varphi) | =  | {G}_{2} |       \\
                                                                               &\implies \frac{ | {G}_{1} |  }{  | \text{ker}(\varphi)  |  }  = | {G}_{2} |  \\ 
                                                                               &\implies | \text{ker}(\varphi)   | = \frac{ | {G}_{1} |  }{ | {G}_{2} |  }.  
            \end{align*}
            Since \( | {G}_{1} |  = 39 \) and \( | {G}_{2} |  = 13 \), then \( | \text{ker}(\varphi)  |  = 3 \). This tells us that \( \text{ker}(\varphi) = \langle \overline{13}   \rangle  \).  
            \end{proof}
        \item[(b)] Suppose that \( {G}_{1} = {D}_{8}  \) and \( {G}_{2} = \Z / 2 \Z \). Determine the possible kernels of \( \varphi \). 
            \begin{proof}
            Again, since \( \varphi  \) is a group homomorphism, the first isomorphism theorem tells us that 
            \[  \frac{ | {G}_{1} |  }{ \text{ker}(\varphi)  }  = 2 \implies | \text{ker}(\varphi)  | = \frac{ 8 }{ 2 }  = 4.   \]
            This tells us that we need to look for all possible subgroups of \( {D}_{8} \) of order \( 4 \). From the lattice of \( {D}_{8} \), we can see that the possible subgroups that \( \text{ker}(\varphi)  \) can be are \( \langle r \rangle \), \( \langle  s r^{2}  \rangle \), and \( \langle rs , r^{2} \rangle \).  
            \end{proof}
        \item[(c)] Suppose that \( | {G}_{1} |   \) and \( | {G}_{2}  |  \) are relatively prime. Prove that \( {G}_{2} = \{ {1}_{G_2} \}  \). 
            \begin{proof}
            Using the first isomorphism theorem and the fact that \( \varphi  \) is a surjective group homomorphism, it follows that   
            \[ \frac{ | {G}_{1} |  }{ | \text{ker}(\varphi)  |  }  = | {G}_{2} |  \implies | {G}_{1} |  = | {G}_{2} | | \text{ker}(\varphi)  |. \]
            Since \( (| {G}_{1} | , | {G}_{2} | ) = 1 \), it follows that \( | {G}_{2} |  = 1  \) and so \( {G}_{2} = \{ {1}_{{G}_{2}} \}  \).
            \end{proof}
    \end{enumerate}
\end{problem}
\begin{problem}
    Let \( G  \) be a group with maximal subgroup \( H \). Recall that \( H  \) is maximal means that \( H \neq G  \) and the only subgroups which contain \( H \) are \( H \) and itself.
    \begin{enumerate}
        \item[(a)] Prove that if \( | G : H  |  = p  \) where \( p  \) is prime, then \( H \) is a maximal subgroup of \( G  \).  
            \begin{proof}
            Our goal is to show that \( H  \) is a maximal subgroup of \( G  \); that is, we need to show that for all \( K \leq G  \) such that \( H \leq K \leq G  \), either \( H = K   \) and \( K = G  \). To this end, let \( K \leq G  \) be such that \( H \leq K \leq G  \). By the Tower Laws, it follows that   
            \begin{align*}
                &| G : H  |  = | G : K  |  | K : H  | \\ 
                &\implies p = | G : K  |  | K  : H |. 
            \end{align*}
            Since \( p \) is prime, either \( | G : K  |  = 1     \) or \( | K : H  |  = 1  \) since the above shows that either \( | G :  K  |  \) or \( | K :  H |  \) divides \( p  \). In either case, \( G = K   \) or \( K =  H \). This tells us that \( H  \) must be maximal.
            \end{proof}
        \item[(b)] Prove that if \( H \trianglelefteq G  \), then \( G / H  \) is cyclic. Deduce that \( | G : H  |  \) is prime.
            \begin{proof}
            Suppose that \( H \trianglelefteq G  \). First, notice that since \( H \trianglelefteq G  \), the group structure of \( G / H  \) is well-defined. Since \( H  \) is maximal, it follows that \( H \neq G  \). Thus, take an element \( g \notin H  \). Since \( g \notin H  \), it follows that \( g H \neq 1 H  \). Consider the cyclic subgroup of \( G / H  \) generated by \( gH  \) denoted by \( \langle gH \rangle \). Since \( H \) is maximal, we know that either \( \langle g H \rangle = G / H   \) or  \( \langle gH  \rangle = \overline{1}  \). Since \( gH \neq 1H  \), we know that \( \langle gH  \rangle = G / H  \) and so \( G / H  \) must be cyclic. Consider the index of \( H  \) in \( G  \) denoted by \( | G : H  |  \). 

            We need to show that \( | G : H |  \) is prime. Suppose for contradiction that \( | G : H  |   \) is composite. Then \( | G : H  | = p q     \) where \( p  \) and \( q  \) are primes. Since \( q  \) divides \( | G : H  |  \), it follows that there exists a subgroup \( K  / H  \) of \( G / H \) of order \( q  \). By the Fourth Isomorphism Theorem, we see that \( K \leq G  \). Since \( H  \) is maximal, \( K = H   \). That is, \( K / H =  \overline{1}  \), implying that \( | K : H  |  = 1  \) which contradics the assumption that \( | K / H  |   \) is prime \( q  \). Thus, \( | G : H  |  = p   \) and so \( | G : H  |  \) must be prime.         
            \end{proof}
    \end{enumerate}
\end{problem}

\begin{problem}
   \begin{enumerate}
       \item[(a)] Suppose \( G  \) is a group with order \( p^{3} q  \) where \( p  \) and \( q  \) are primes with \( p > q  \). Prove that if \( G  \) has a subgroup \( H  \) of order \( p^{3} \), then it must be normal. 
           \begin{proof}
           Consider the quotient \( G  / H \). Then its order is     
           \[ | G / H  |  = \frac{ | G |  }{ | H |  } = \frac{ p^{3} q  }{ p^{3} } = q \implies G /H \ \text{is cyclic.}    \]
           This tells us further that \( G / H  \) is abelian, and so every subgroup of \( G / H \) is normal. By the fourth isomorphism theorem, we know that \( H / H   \trianglelefteq G / H \) implies that \( H \trianglelefteq G   \).  
           \end{proof}
       \item[(b)] Prove that if \( H \leq Z(G) \), then \( G  \) must be abelian.  
           \begin{proof}
           Notice that the cosets of \( H  \) in \( G  \) partition \( G  \); that is,
           \[  H = \bigcup_{ g \in G  }^{  }  gH. \]
           Now, let \( x,y \in G \). Then for some \( {g}_{1}, {g}_{2} \in G  \), we have \( x \in {g}_{1} H  \) and \( y \in {g}_{2} H  \). That is, for some \( {h}_{1}, {h}_{2} \in H  \), we have \( x = {g}_{1} {h}_{1} \) and \( y = {g}_{2} {h}_{2} \). Now, we have since \( {h}_{1}, {h}_{2} \in H \) and \( H \leq Z(G) \), we get that \( {h}_{1}, {h}_{2} \in Z(G) \) and so  
           \begin{align*}
               xy &= ({g}_{1}{h}_{1})({g}_{2} {h}_{2})  \\
                  &= {g}_{1}({g}_{2}{h}_{1}){h}_{2} \\
                  &= ({h}_{1}{h}_{2})({g}_{1}{g}_{2}) \\
                  &= ({g}_{2}{h}_{2})({g}_{1}{h}_{1}) \\
                  &= yx.
           \end{align*}
           Thus, \( xy = yx  \) and so \(  G  \) must be abelian. 
           \end{proof}
   \end{enumerate}  
\end{problem}
\vspace{1mm}
\begin{problem} Let \( G  \) be a group and \( g \in G  \). Define \( {\sigma}_{g} (x) = g x g^{-1} \) for all \( x \in G  \).
   \begin{enumerate}
       \item[(a)] Prove that the mapping \( \psi : G \to {S}_{G} \) defined by \( \psi(g) = {\sigma}_{g} \) is a homomorphism.
           \begin{proof}
           For all \( {g}_{1}, {g}_{2} \in G  \) and for all \( x \in G  \), we have 
           \begin{align*}
               \psi({g}_{1}{g}_{2})(x) = {\sigma}_{{g}_{1}{g}_{2}}(x) &= ({g}_{1}{g}_{2})x ({g}_{1} {g}_{2})^{-1}  \\
                                                                      &={g}_{1} ({g}_{2} x {g}_{2}^{-1}) {g}_{1}^{-1} \\
                                                                      &= {g}_{1}({\sigma}_{{g}_{2}}(x)) {g}_{1}^{-1} \\
                                                                      &= {\sigma}_{{g}_{1}}({\sigma}_{{g}_{2}}(x)) \\
                                                                      &= (\sigma_{{g}_{1}} \circ \sigma_{{g}_{2}})(x) \\
                                                                      &= (\psi({g}_{1}) \circ \psi({g}_{2})) (x).
           \end{align*}
           This implies that \( \psi({g}_{1}{g}_{2}) = \psi({g}_{1})\circ \psi({g}_{2}) \) and so we conclude that \( \psi \) is a group homomorphism.
           \end{proof}
        \item[(b)] Prove that \( G / Z(G) \) is isomorphic to a subgroup of \( {S}_{G} \).
            \begin{proof}
            Using part (a), it follows that 
            \[  G / \text{ker}(\psi) \cong \psi(G) \leq {S}_{G}.  \]
            From here, it suffices to show that \( \text{ker}(\psi) = Z(G)  \). Indeed, we get that 
            \begin{align*}
            \text{ker}(\psi) &= \{ g \in G : \psi(g)(x) = x \ \forall x \in G    \}  \\
                             &= \{ g \in G : \sigma_g(x) = x \ \forall x \in G  \}  \\
                             &= \{  g \in G : g x g^{-1}  =x \ \forall  x \in G  \} \\
                             &= Z(G).
            \end{align*}
            Thus, we have 
            \[  G / Z(G) \cong K \leq {S}_{G}  \]
            where \( K = \psi(G) \).
            \end{proof}
   \end{enumerate} 
\end{problem}

\end{document}

